\documentclass[a4paper,12pt]{article}

% --- Use XeLaTeX font support ---

\usepackage{fontspec}
\setmainfont{TeX Gyre Termes} % nice Times-like font, installed by texlive-core
\setsansfont{TeX Gyre Heros} % optional sans-serif
\setmonofont{Fira Code}       % optional monospaced font
% --- Math packages ---
\usepackage{amsmath, amssymb, amsthm}
\usepackage{mathtools}

% --- Page layout ---
\usepackage[margin=1in]{geometry}

% --- Better lists ---
\usepackage{enumitem}
% --- Hyperlinks for references ---
\usepackage[colorlinks=true, linkcolor=blue, urlcolor=blue]{hyperref}

% Header: fancy layout like the rovided image
\usepackage{fancyhdr}
\setlength{\headheight}{26pt} % increase if fancyhdr warns
\pagestyle{fancy}
\fancyhf{} % clear header/footer
\rhead{\begin{tabular}{r} Mt.: 3052737\end{tabular}}
\chead{\begin{tabular}{c}Lösungen EA 2\end{tabular}}
\lhead{\begin{tabular}{l}61521 Einführung in die \\ Numerische Mathematik SS26 \end{tabular}}
\renewcommand{\headrulewidth}{0.8pt}
\fancyfoot[C]{- \thepage\ -}
% Ensure the plain page style (used by \maketitle) uses the same header

% --- Line spacing ---
\usepackage{setspace} % provides \doublespacing and \onehalfspacing

% --- Optional solutions environment ---
\newenvironment{solution}{
  \par\noindent\textbf{Solution:}\quad
}{\hfill$\blacksquare$\par\vspace{1em}}

% % --- Title info ---
% \title{Aufgabe \#1}
% \author{Your Name}
% \date{\today}
% 
\begin{document}
\raggedright
% Enable double (2.0) line spacing for the document body
\doublespacing
\noindent\textbf{Aufgabe 2} \\
\noindent\textbf{(a)} Sei der gesamte Grenzwert $x$. Der Ausdruck lässt sich dan schreiben als 
$$x= \sqrt{3+2⋅(Rest)}$$
Der Rest ist genau derselbe unendliche Ausdruck wie das Ganze, also $x= \sqrt{3+2x}$.
Die zugehörige Iterationsfunktion lautet somit: 
$$\varphi: \mathbb{R} \rightarrow \mathbb{R}, \ \varphi(x) = \sqrt{3 + 2x}.$$
Der Fixpunkt lässt sich leicht als $3$ ermitteln.

Nachweis der Fixpunkt-Eigenschaft:

\begin{enumerate}
    \item Fixpunkt bei $x = 3$:
    \[
    \varphi(3) = \sqrt{3 + 2 \cdot 3} = \sqrt{9} = 3 
    \]

    \item Definitionsbereich:
    \[
    \varphi : \left[-\frac{3}{2}, \infty\right) \to [0, \infty)
    \]
    (Damit die Wurzel definiert ist)

    \item Konvergenz (für Startwerte $x_0 \geq 0$):
    Die Ableitung ist:
    \[
    \varphi'(x) = \frac{2}{2\sqrt{3 + 2x}} = \frac{1}{\sqrt{3 + 2x}}
    \]
    Am Fixpunkt:
    \[
    |\varphi'(3)| = \frac{1}{\sqrt{9}} = \frac{1}{3} < 1
    \]
    Nach dem Banachschen Fixpunktsatz konvergiert die Iteration
    \[
    x_{n+1} = \varphi(x_n)
    \]
    für alle Startwerte $x_0 \geq 0$ gegen den Fixpunkt $x^* = 3$.
\end{enumerate}

Antwort: $\varphi(x) = \sqrt{3 + 2x}$ mit Startwert $x_0 \geq 0$.

\noindent\textbf{(b)} Prüfen der Vorraussetzungen des Satzes 2.2.1:

\begin{itemize}
\item Abgeschlossenheit: Wir wählen das abgeschlossene Intervall $I = [0, 4].$
\item Selbstabbildung ($\varphi(I) \subseteq I$):
Die Funktion $\varphi(x) = \sqrt{3 + 2x}$ ist auf $I$ streng monoton wachsend.
Wir prüfen die Randpunkte:\\
$$\varphi(0) = \sqrt{3 + 0} = \sqrt{3} \approx 1{,}732 \ge 0$$ 
$$\varphi(4) = \sqrt{3 + 8} = \sqrt{11} \approx 3{,}317 \le 4$$\\
Daraus folgt, dass $\varphi([0, 4]) = [\sqrt{3}, \sqrt{11}] \subset [0, 4]$ ist. Die Funktion bildet das Intervall also in sich selbst ab. 
\item Kontraktionseigenschaft ($|\varphi'(x)| \le L < 1$):
$$\varphi'(x) =  \frac{1}{\sqrt{3 + 2x}}$$
Da der Nenner für größer werdende $x$ wächst, ist $\varphi'(x)$ streng monoton fallend auf $I$. Ihr Maximum nimmt die Ableitung also am linken Rand $x = 0$ an:$$\max_{x \in [0, 4]} |\varphi'(x)| = \varphi'(0) = \frac{1}{\sqrt{3}} \approx 0{,}577<1$$
Es handelt sich also um eine Kontraktion.

\end{itemize}
Da alle Voraussetzungen erfüllt sind, existiert auf dem Intervall $I = [0, 4]$ nach dem Banachschen Fixpunktsatz genau ein Fixpunkt und die Iteration konvergiert für jeden Startwert $x_0 \in I$ gegen diesen.\\


\bigskip
\noindent\textbf{(c)} Berechnung der ersten 5 Iterierten: \\
\begin{align*}
x_0 &= \sqrt{3} \approx 1{,}73205\\
x_1 &= \sqrt{3 + 2 \cdot 1{,}73205} = \sqrt{6{,}46410} \approx 2{,}54246\\
x_2 &= \sqrt{3 + 2 \cdot 2{,}54246} = \sqrt{8{,}08492} \approx 2{,}84340\\
x_3 &= \sqrt{3 + 2 \cdot 2{,}84340} = \sqrt{8{,}68680} \approx 2{,}94734\\
x_4 &= \sqrt{3 + 2 \cdot 2{,}94734} = \sqrt{8{,}89468} \approx 2{,}98239\\
x_5 &= \sqrt{3 + 2 \cdot 2{,}98239} = \sqrt{8{,}96478} \approx 2{,}99413 
\end{align*}
Die Formel für die a-posteriori-Fehlerabschätzung lautet (Satz 2.2.2):
$$|x_n - x^*| \le \frac{L}{1-L} |x_n - x_{n-1}|$$
Aus (b):
$$\frac{L}{1-L} = \frac{\frac{1}{\sqrt{3}}}{1 - \frac{1}{\sqrt{3}}} = \frac{1}{\sqrt{3}-1} \approx 1{,}36603$$
Differenz der letzten beiden Iterierten:$$|x_5 - x_4| \approx |2{,}99413 - 2{,}98239| = 0{,}01174$$
Eingesetzt in die Abschätzungsformel ergibt sich die Fehlerschranke:$$|x_5 - x^*| \le 1{,}36603 \cdot 0{,}01174 \approx 0{,}01604$$Der Fehler für $x_5$ liegt laut Abschätzung also garantiert unter $0{,}01604$
\end{document}
